% !TEX root = ../../../main.tex
% 来源：高中物理甲种本·第三册（电磁·光学·近代物理）
% 章节：交流电 / 滤波
% 原始文件：3.tex，行 1469-1491
% 类型：circuit
% ---
\begin{circuitikz}[>=latex, scale=1.3]

\draw (-3,.25)--(-3,1)--(-2+.25,1) to [american inductor] (-2+.25,-1)--(-3,-1)--(-3,-.25);
\draw (1,-1)--(1,-1.5)--(-1.25,-1.5) to [american inductor] (-1.25,1.5)--(1,1.5)--(1,1);
\draw [ultra thick] (-3/2, -.4)--(-3/2, .4);
\ctikzset{diodes/scale=0.6} \draw (0,0) to [full diode] (1,1) to [full diode] (2,0);
\ctikzset{diodes/scale=0.6} \draw (0,0) to [full diode] (1,-1) to [full diode] (2,0) ;
\draw (2,0)to [american inductor](5,0)--(5.5,0) to [european, R=$R$] (5.5,-2) --(0,-2)--(0,0);
\draw [ultra thick] (3+.1, .25) to node [above]{$L$} (4-.1, .25);


\draw  (2.5,0) to [eC=$C_1$, *-*] (2.5,-2);
\draw (4.2,0) to [eC=$C_2$, *-*] (4.2,-2);

\draw [fill=white] (-3,.25) circle (1pt);
\draw [fill=white] (-3,-.25) circle (1pt);
\draw [fill=black] (0,0) circle (1pt);
\draw [fill=black] (2,0) circle (1pt);\draw [fill=black] (1,-1) circle (1pt);
\draw [fill=black] (1,1) circle (1pt);
\node at (-3,0){\Huge $\sim $};


\end{circuitikz}
